% !PRINTED 250
\weberpage{250}
%<<CONT
first term is a square and whose last is a non-square, and choose in it a square term followed by a non-square.

If then $a = \xi^2$ is a square and $\xi^2 + 1$ a non-square, then, if $\beta$ is a non-square, $\beta : \xi^2 + 1 = \eta^2$ is a square, and from this it follows:
\[ \beta = \xi^2 \eta^2 + \eta^2 \]
q. e. d.

\begin{center}
\S.\,65.

\gesperrt{Congruence groups in the field $\mathfrak{C}$.}
\end{center}

In every congruence field $\mathfrak{C}$ one can form a finite group of linear substitutions, if in
\weq{1}{A = \begin{pmatrix} \alpha, & \beta \\ \gamma, & \delta \end{pmatrix}}
one lets the elements $\alpha, \beta, \gamma, \delta$ run through all the numbers of $\mathfrak{C}$, and if one composes any two of these substitutions according to the rule
\weq{2}{\begin{pmatrix} \alpha, & \beta \\ \gamma, & \delta \end{pmatrix} \begin{pmatrix} \alpha', & \beta' \\ \gamma', & \delta' \end{pmatrix} = \begin{pmatrix} \alpha\alpha' + \beta\gamma', & \alpha\beta' + \beta\delta' \\ \gamma\alpha' + \delta\gamma', & \gamma\beta' + \delta\delta' \end{pmatrix}}
The degree of this group is $p^{4n}$, if $p$ is the modulus and $n$ the degree of $\mathfrak{C}$.

But in it there is contained a group of lower degree. If we call $\alpha\delta - \beta\gamma$ the determinant of the substitution (1), it follows from (2) that the determinant of the substitution $AA'$ composed from $A$ and $A'$ is equal to the product of the determinants of $A$ and $A'$. If therefore we impose on the numbers $\alpha, \beta, \gamma, \delta$ the condition that
\weq{3}{\alpha\delta - \beta\gamma = 1}
shall hold, then these elements $A$ likewise form a group. The degree of this group is equal to the number of solutions of (3). To find this number, we first observe that for $\alpha, \beta$ we may take any two numbers from $\mathfrak{C}$ which are not both equal to zero. For if, say, $\alpha$ is different from zero, one can set $\gamma = 0, \delta = 1 : \alpha$ and thereby satisfy the condition (3). The number of usable pairs of numbers $\alpha, \beta$ is therefore $p^{2n} - 1$. But if for a pair of numbers $\alpha, \beta$ a solution $\gamma_0, \delta_0$ of (3) has been found, then all the others must satisfy the condition
\[ \alpha (\delta - \delta_0) = \beta (\gamma - \gamma_0) \]
%CONT>>
